



We provide an overview of our main results here.
We begin with a brief introduction to the necessary notation and background for providing a clear overview.
The full preliminaries can be found in \Cref{sec:preliminaries}.



We use superscripts to indicate the systems on which linear operators and linear maps are defined.
For instance, $M^{\sA}$ denotes a linear operator on system $\sA$ and $\cT^{\sA \rarr \sB}$ denotes a linear map from $\sA$ to $\sB$, and additionally, $\cT^{\sA \rightarrow \sA}$ is simply written as $\cT^{\sA}$.
The superscript is omitted when it is clear from the context.
The identity operator $\bI$ and the identity map $\rid$ are often implicit; for instance, we denote $\bI^\sA \otimes M^\sB$ simply by $M^{\sB}$.
For quantum states, we denote a pure state on $\sA\sB$ such as $\ket{\omega}^{\sA\sB}$, and its reduced density operator on $\sA$ by $\omega^\sA$, i.e., $\omega^\sA = \tr_\sB\big[\ketbra{\omega}{\omega}^{\sA\sB}\big]$, where $\tr_\sB$ is the partial trace over the system $\sB$.

The (squared) fidelity between quantum states $\rho$ and $\sigma$ is defined as $\rF(\rho, \sigma) = \big(\tr\big[\sqrt{\sqrt{\sigma}\rho\sqrt{\sigma}}\big]\big)^2$.
We also use the square root fidelity, defined as $\sqrt{\rF}(\rho, \sigma) = \sqrt{\rF(\rho, \sigma)}$.
The fidelity is monotonic with respect to the partial trace: $\rF(\rho^\rA, \sigma^\rA) \geq \rF(\rho^{\rA\rB}, \sigma^{\rA\rB})$.
For notational convenience, we write the fidelity as $\rF(\rho, \ket{\phi})$ or $\rF(\ket{\psi}, \ket{\phi})$, when one or both of its arguments are pure states.





Our goal is to construct a quantum algorithm that \emph{approximately} realizes the unitary $U^\sB$ achieving Eq.~\eqref{inteq:57}.
In the algorithm construction, we consider the following three computational models:
\begin{enumerate}[label=\Roman*.]
    \item Purified query access model~\cite{watrous2002limitQSZK, watrous2009QSZKquantumattacks, belovs2019ClassicalDist, gilyen2019DistPropTesting, brandao2019SDPOpt, vanapeldoorn2019impleveSDP, Subramanian2021estimatealphalenni, gilyen2022improvedfidelity, Apeldoorn2023TomoPurifQuery, wang2023QAfidelityest, wang2023SampletoQuary, luo2024SuccinctTest, wang2024NewQuantEnt, Liu2024geometricmean}: let $U_\rho^{\sA\sB}$ and $U_\sigma^{\sA\sB}$ be unitaries which prepare states $\ket{\rho}^{\sA\sB}$ and $\ket{\sigma}^{\sA\sB}$; respectively, that is, $U_\rho^{\sA\sB}\ket{0}^{\sA\sB}=\ket{\rho}^{\sA\sB}$, and $U_\sigma^{\sA\sB}\ket{0}^{\sA\sB}=\ket{\sigma}^{\sA\sB}$.
    We assume that $U_\rho^{\sA\sB}$ and $U_\sigma^{\sA\sB}$ are available multiple times as unitary oracles.
    It is also assumed that we can query their inverses $(U_\rho^{\sA\sB})^\dag$ and $(U_\sigma^{\sA\sB})^\dag$.
    \item Purified sample access model~\cite{chen2024localtestUniInv, liu2024exponentialseparations}:
    multiple independent and identical copies of purified states $\ket{\rho}^{\sA\sB}$ and $\ket{\sigma}^{\sA\sB}$ are available.
    \item Mixed sample access model~\cite{gilyen2022improvedfidelity, wang2023SampletoQuary, wang2024TimeEfficientEntEstSamp, wang2024sampleoptimal}: multiple independent and identical copies of states $\rho^\sA$ and $\sigma^\sA$ are available.
\end{enumerate}

For each model, we evaluate the number of queries/samples for implementing the Uhlmann transformation.
Since each model listed above can simulate those listed below, the models are ordered from the one with the strongest assumptions and highest computational power (I) to the one with the weakest assumptions and lowest power (III).
The details and motivations behind these computational models are presented in \Cref{sec:three computational models}.

When we consider the Uhlmann transformation, the purified sample access model (II) appears to be the most natural setting, as the transformation is specified for a pair of purified states and cannot be determined from mixed states alone.
Nevertheless, there is still motivation to consider the mixed sample access model (III), due to its utility in various information-theoretic tasks, as discussed in our applications.
In the mixed sample access model (III), we consider the Uhlmann transformation tailored to a specific purification named the \emph{canonical purification} (its definition is given by Eq.~\eqref{eq:def of canonical purification}).
We construct a quantum algorithm that prepares the canonical purification from multiple sampled states, and utilize it as part of our Uhlmann transformation algorithm in model (III).


We further introduce some technical notation used throughout this paper; see \Cref{tab:technical notation} for a summary.
These quantities are related to each other. For example, $r_\omega \leq \kappa_\omega$ holds for any state $\omega$, while $r_\omega$ does not imply any upper bound on $\kappa_\omega$.
The inequality $r \leq r_{\rm min}$ follows from the fact that the rank of a product of matrices is at most the rank of each matrix.
Moreover, since $r s_{\rm min} \leq \rF(\rho^\sA, \sigma^\sA) \leq 1$, it follows that $r \leq 1/s_{\rm min}$.





%===================================================================================================================================================================================================================



\renewcommand{\arraystretch}{1.8}
\begin{table*}
\centering
\caption{A table of some technical notation.}
\begin{tabular}{wc{5em}|wl{33em}} 
    $d_\sA$, $d_\sB$ & \hspace{1mm} The dimensions of the Hilbert spaces $\cH^\sA$ and $\cH^\sB$, respectively. \\ \hline
    $s_{\rm min}$, $r$ & \hspace{1mm} The minimum non-zero singular value and the rank of $\sqrt{\sigma^\sA}\sqrt{\rho^\sA}$, respectively. \\ \hline
        $r_\rho$, $r_\sigma$& \hspace{1mm} The ranks of $\rho^\sA$ and $\sigma^\sA$, respectively. \\ \hline
    $r_{\rm min}$ & \hspace{1mm} The smaller of the ranks of $\rho^\sA$ and $\sigma^\sA$: $r_{\rm min} = \min\{r_\rho, r_\sigma\}$. \\ \hline 
    $\rho_{\rm min}$, $\sigma_{\rm min}$ & \hspace{1mm} The minimum non-zero eigenvalues of $\rho^\sA$ and $\sigma^\sA$, respectively. \\ \hline
    $\kappa_\rho$, $\kappa_\sigma$ & \hspace{1mm} The reciprocals of $\rho_{\rm min}$ and $\sigma_{\rm min}$, respectively: $\kappa_\rho = 1/\rho_{\rm min}$ and $\kappa_\sigma = 1/\sigma_{\rm min}$. \\
    \hline
    $\kappa_{\rm min}$ & \hspace{1mm} The smaller of the reciprocals of $\rho_{\rm min}$ and $\sigma_{\rm min}$: $\kappa_{\rm min} = \min\{\kappa_\rho, \kappa_\sigma\}$.
\end{tabular}
\label{tab:technical notation}
\end{table*}
\renewcommand{\arraystretch}{1.0}



%===================================================================================================================================================================================================================



For evaluating computational costs, we use the Landau notation $\cO(\cdot)$, $\Omega(\cdot)$, and $\til{\cO}(\cdot)$. These describe the scaling in terms of dimension $d$ and the approximation accuracy $\delta$, that is, a notation such as $\cO(d/\delta)$ is understood as both $\cO(d)$ and $\cO(1/\delta)$.
In particular, $\til{\cO}\big(f(d, 1/\delta)\big)$ represents $\til{\cO}\big(f(d, 1/\delta)\big) = \cO\big(f(d, 1/\delta) {\rm polylog}(d, 1/\delta)\big)$,
where $f$ is a certain function; it hides polylogarithmic factors in $d$ and $1/\delta$. 
We also use the notation $\poly(d, 1/\delta)$ to denote a polynomial in $d$ and $1/\delta$.





We first provide upper bounds on the computational costs for implementing the Uhlmann transformation in \Cref{sec:upper bound uhl}.
Next, in \Cref{sec:comparison with state tom}, we compare the computational costs of our algorithm with other algorithms, highlighting the advantage of our algorithm.
In \Cref{sec:overview of lower bound}, we provide a lower bound on the query and sample complexities for generally implementing the Uhlmann transformation, which follows from the hardness of the fidelity estimation.
In \Cref{sec:overview application}, we present applications of the Uhlmann transformation algorithm.









%========================================================================================================================================================================================================================================================================================================







\subsection{Upper bounds on the computational costs via Uhlmann transformation algorithms}
\label{sec:upper bound uhl}



We here provide upper bounds on the computational costs required for the Uhlmann transformation in each computational model in Secs.~\ref{sec:upper purif query main},~\ref{sec:upper purif sample main}, and~\ref{sec:upper mixed sample main}.





\subsubsection{In the purified query access model}
\label{sec:upper purif query main}


Our statement in the purified query access model is as follows.


\begin{theorem}[Uhlmann transformation algorithm in the purified query access models]
\label{infthm:uhlmenn purif query}
For any $\delta \in (0, 1)$, there exists a quantum query algorithm that realizes a quantum channel $\cT^{\sB}$ satisfying 
\begin{align}
\label{inteq:17}
     \rF(\cT^{\sB}(\ketbra{\rho}{\rho}^{\sA\sB}), \ket{\sigma}^{\sA\sB}) \geq \rF(\rho^\sA, \sigma^\sA) - \delta,
\end{align}
using $u$ queries to $U_\rho^{\sA\sB}$, $U_\sigma^{\sA\sB}$, and their inverses, where $u = \cO\big(\min\big\{\f{1}{s_{\rm min}}, \f{r}{\delta}\big\}\log{\big(\f{1}{\delta}\big)}\big)$. The quantum circuit of the algorithm consists of $\cO\big(u\log{(d_\sA d_\sB)}\big)$ one- and two-qubit gates, and at any one time, $\cO\big(\log(d_\sA d_\sB)\big)$ qubits suffice.
\end{theorem}


A more formal and technical version of \Cref{infthm:uhlmenn purif query} is provided in \Cref{sec:purif query} as \Cref{thm:Uhlmann alg purif query model}, where the approximation error of the transformation is measured not only by the fidelity difference but also by the diamond norm, in order to evaluate the performance for quantum channels.



Theorem~\ref{infthm:uhlmenn purif query} states that, by this quantum query algorithm, the Uhlmann transformation is approximately realized with accuracy $\delta$.
Remarkably, when either $1/s_{\rm min}$ or $r/\delta$ is sufficiently small, for instance, a polynomial in the number of qubits, the algorithm is efficiently implementable with a polynomial number of queries.
This implies that our approach achieves an exponential improvement in query and gate costs over the previous approach in Ref.~\cite{metger2023stateqipspace}, which requires exponential costs, $\poly(d_\sA d_\sB, 1/\delta)$.
We will discuss comparison in detail in \Cref{sec:comparison with state tom}.




This algorithm is a query algorithm. However, if quantum circuit descriptions of the unitaries $U_\rho^{\sA\sB}$ and $U_\sigma^{\sA\sB}$ are given, in other words, if one can implement $U_\rho^{\sA\sB}$, $U_\sigma^{\sA\sB}$, their inverses on one's own, the circuit complexity, the total number of one- and two-qubit gates required in the algorithm, is evaluated as
\begin{align}
    \cO\Big(u \big(\cC(U_\rho) + \cC(U_\sigma) + \log{(d_\sA d_\sB)}\big)\Big),
\end{align}
where $\cC(U_\rho)$ and $\cC(U_\sigma)$ denote the circuit complexity of $U_\rho^{\sA\sB}$ and $U_\sigma^{\sA\sB}$, respectively.


One might be concerned about the necessity of knowing either the value of $s_{\rm min}$ or $r$ for implementing this algorithm.
To estimate the singular value, quantum algorithms based on the quantum phase estimation have been proposed in Refs.~\cite{chakraborty2019powerblock, Gilyn2019thesis}.
We can estimate $s_{\rm min}$ with the desired accuracy in advance.
Similarly, several algorithms for rank estimation are known~\cite{Tan2021ValEstRank, wang2024NewQuantEnt}.
In our construction of the Uhlmann transformation algorithm, we employ the QSVT with the sign function introduced in \Cref{sec:BE and QSVT intro}.
Due to a property of the QSVT with the sign function, it is sufficient to have a lower bound on $s_{\rm min}$ or an upper bound on $r$.









\subsubsection{In the purified sample access model}
\label{sec:upper purif sample main}


Next, we provide the statement in the purified sample access model.

\begin{theorem}[Uhlmann transformation algorithm in the purified sample access models]
\label{infthm:uhlmenn purif sample}
For any $\delta \in (0, 1)$, there exists a quantum sample algorithm that realizes a quantum channel $\cT^{\sB}$ satisfying
\begin{align}
\label{eq:fid diff pure samp}
     \rF(\cT^{\sB}(\ketbra{\rho}{\rho}^{\sA\sB}), \ket{\sigma}^{\sA\sB}) \geq \rF(\rho^\sA, \sigma^\sA) - \delta,
\end{align}
using $w$ samples of $\ket{\rho}^{\sA\sB}$ and $\ket{\sigma}^{\sA\sB}$, where $w = \cO\big(\f{1}{\delta}\min\big\{\f{1}{s_{\rm min}^2}, \f{r^2}{\delta^2}\big\}\big(\log{\big(\f{1}{\delta}\big)}\big)^2\big)$. The quantum circuit of this algorithm consists of $\cO\big(w\log (d_\sA d_\sB)\big)$ one- and two-qubit gates, and at any one time, $\cO\big(\log{(d_\sA d_\sB)}\big)$ qubits suffice.
\end{theorem}

A more formal and technical version of \Cref{infthm:uhlmenn purif sample} is \Cref{thm:algorithm Uhlmann} in \Cref{sec:purif sample}, where the approximation error is evaluated by both fidelity difference and diamond norm.

Theorem~\ref{infthm:uhlmenn purif sample} states that, rather remarkably, the sample complexity still does not scale with the dimension $d_\sA$ or $ d_\sB$---when either $1/s_{\rm min}^2$ or $r^2/\delta^2$ is sufficiently small, the Uhlmann transformation can be efficiently implemented using multiple copies of the purified states $\ket{\rho}^{\sA\sB}$ and $\ket{\sigma}^{\sA\sB}$. 
Moreover, because this algorithm is sequential and processes the quantum states one by one, $\cO\big(\log(d_\sA d_\sB)\big)$ qubits suffice at any one time, which includes qubits to store these states.






%=====================================================================================================================================================================================================







%=====================================================================================================================================================================================================




\subsubsection{In the mixed sample access model}
\label{sec:upper mixed sample main}




In general, purifications of a mixed state are not unique. 
This implies that the Uhlmann transformation is not determined solely by a given pair of mixed states.
In other words, in the mixed sample access model, the Uhlmann transformation is ill-defined in the sense that, from samples of mixed states, it is impossible to perform the Uhlmann transformation that universally works for all possible purifications.
We thus first perform a specific purification using multiple copies of the given mixed states.
This allows us to implement the Uhlmann transformation between the specific purified states obtained by this purification.

In this work, we consider the canonical purification: for any state $\omega^\sA$, the canonical purified state $\ket{\omega_{\rm c}}^{\sA\sB}$ is defined by 
\begin{equation}
\label{eq:def of canonical purification}
    \ket{\omega_{\rm c}}^{\sA\sB} = \big(\sqrt{\omega^\sA} \otimes \bI^{\sB}\big)\ket{\Gamma}^{\sA\sB},
\end{equation}
where $\ket{\Gamma}^{\sA\sB} = \sum_{i=1}^{d_\sA}\ket{i}^{\sA}\ket{i}^{\sB}$ is the unnormalized maximally entangled state in the computational basis $\ket{i}$. 
It is natural to assume that $d_\sA = d_\sB$ here.
The state $\ket{\omega_{\rm c}}^{\sA\sB}$ is indeed one of the purified states of $\omega^\sA$ as $\tr_\sB[\ketbra{\omega_{\rm c}}{\omega_{\rm c}}^{\sA\sB}] = \omega^\sA$.
A quantum algorithm for the canonical purification from multiple identical copies of $\omega^\sA$ is provided in \Cref{sec:appendix B}.


Let $\ket{\rho_{\rm c}}^{\sA\sB}$ and $\ket{\sigma_{\rm c}}^{\sA\sB}$ be the canonical purified states of $\rho^\sA$ and $\sigma^\sA$, respectively.
Our statement for the mixed sample access model is as follows.


\begin{theorem}[Uhlmann transformation algorithm in the mixed sample access models]
\label{infthm:uhlmenn mixed sample}
For any $\delta \in (0, 1)$, there exists a quantum sample algorithm that realizes a quantum channel $\cT^\sB$ satisfying 
\begin{align} 
    \rF\big(\cT^\sB(\ketbra{\rho_{\rm c}}{\rho_{\rm c}}^{\sA\sB}), \ket{\sigma_{\rm c}}^{\sA\sB}\big)  
    \geq \rF(\rho^\sA, \sigma^\sA) - \delta, 
\end{align}
using $\zeta$ samples of $\rho^\sA$ and $\sigma^\sA$, where $\zeta = \til{\cO}\Big(\f{d_\sA}{\delta^2\beta^3}\min\Big\{\f{1}{\rho_{\rm min}^2} + \f{1}{\sigma_{\rm min}^2}, \f{d_\sA^2}{\delta^4\beta^4}\Big\}\Big)$, and $\beta = \cO(\max\{s_{\rm min}, \delta/r\})$. The quantum circuit of this algorithm consists of $\cO(\zeta\log{d_\sA})$ one- and two-qubit gates, and at any one time, $\cO(\log{d_\sA})$ qubits suffice.
\end{theorem}




A formal and technical version of \Cref{infthm:uhlmenn mixed sample} is given by \Cref{thm:Uhlmann mixed state sample} in \Cref{sec:mixed sample}, where the approximation error is evaluated by both fidelity difference and diamond norm.













Due to the symmetry of $\ket{\Gamma}^{\sA\sB}$, the canonical purification has the property that the reduced states on each system are related by the transpose. That is, for the canonical purified state $\ket{\rho_{\rm c}}^{\sA\sB}$ of a state $\rho$ on $\sA$, the reduced states are $\rho_{\rm c}^\sA = \rho$ and $\rho_{\rm c}^\sB = \rho^\top$.
For $\rho_{\rm c}^\sB = \rho^\top$ and $\sigma_{\rm c}^\sB = \sigma^\top$, the Uhlmann unitary $U^\sA$ satisfies $\rF(\rho^\top, \sigma^\top) = \rF(U^\sA\ket{\rho_{\rm c}}^{\sA\sB}, \ket{\sigma_{\rm c}}^{\sA\sB})$, where the states with transpose in the left-hand side are defined on $\sB$.
Since the transpose does not change the fidelity, we obtain $\rF(\rho^\sA, \sigma^\sA) = \rF(U^\sA\ket{\rho_{\rm c}}^{\sA\sB}, \ket{\sigma_{\rm c}}^{\sA\sB})$.
This shows that, when we consider the canonical purification, we can apply a transformation by $U^\sA$ to the original system $\sA$ to achieve the fidelity of the states on the original system. 
This is unlike the original Uhlmann transformation, which is applied to the purifying system $\sB$.
We call the transformation on $\sA$ a \emph{variant} of the Uhlmann transformation and propose an algorithm for implementing it in the mixed sample access model.

Note that such a unitary $U^\sA$ exists due to the property of the canonical purification.
For general purifications, applying a unitary to the original system $\sA$ and achieving the Uhlmann transformation are not achievable.



\begin{theorem}[Algorithm for a variant of the Uhlmann transformation in the mixed sample access models]
\label{thm:variant Uhl overview}
For any $\delta \in (0, 1)$, there exists a quantum sample algorithm that realizes a quantum channel $\cT^\sA$ satisfying 
\begin{align} 
    \rF\big(\cT^\sA(\ketbra{\rho_{\rm c}}{\rho_{\rm c}}^{\sA\sB}), \ket{\sigma_{\rm c}}^{\sA\sB}\big)  
    \geq \rF(\rho^\sA, \sigma^\sA) - \delta, 
\end{align}
using $\zeta$ samples of $\rho^\sA$ and $\sigma^\sA$, where $\zeta = \til{\cO}\big(\f{1}{\delta\beta}\min\big\{\f{1}{\rho_{\rm min}^2} + \f{1}{\sigma_{\rm min}^2}, \f{1}{\beta^4\delta^4}\big\}\big)$, and $\beta = \cO(\max\{s_{\rm min}, \delta/r\})$.
The quantum circuit of this algorithm consists of $\cO(\zeta\log{d_\sA})$ one- and two-qubit gates, and at any one time, $\cO(\log{d_\sA})$ qubits suffice.
\end{theorem}


A formal and technical version of \Cref{thm:variant Uhl overview} is given by \Cref{thm:variant of uhlmann mixed sample} in \Cref{sec:altanative mixed sample uhlmann}, where the approximation error is evaluated using both fidelity difference and the diamond norm.


Compared with the algorithm in \Cref{infthm:uhlmenn mixed sample}, this algorithm in \Cref{thm:variant Uhl overview} can be implemented with substantially fewer samples of the mixed states.
This is because this algorithm does not rely on the quantum algorithm for the canonical purification, which requires a large number of samples. 
Instead, it directly block-encodes the square root of the given mixed states into a large unitary through a combination of the QSVT and the density matrix exponentiation.
Further details on this variant of the Uhlmann transformation can be found in \Cref{sec:altanative mixed sample uhlmann}.







\renewcommand{\arraystretch}{2.4}
\begin{table*}
\centering
\caption{Comparison of the query/sample complexity for the Uhlmann transformation. We here focus only on the results with better scaling with respect to the accuracy $\delta$ for simplicity.
For the notation, see \Cref{tab:technical notation}.
The result of the state tomography-based approach is derived based on the results of, to the best of our knowledge, the most efficient quantum state tomography algorithms~\cite{ODonnell2016EffTomo, Haah2017samploptTomo, Apeldoorn2023TomoPurifQuery}.
Notably, in the purified query access model and the purified sample access model, the query and sample complexities of our algorithms are independent of the dimensions $d_\sA$ and $d_\sB$, which indicates exponential advantages.}
\begin{tabular}{wc{7em}|wc{9em}|wc{10em}|wc{12em}}
                & Our algorithm  & State tomography-based & Prior work \\ \hline
    Purified query &   $\cO\big(\log{(1/\delta)}/{s_{\rm min}}\big)$ & $\til{\cO}\big(d_\sA d_\sB / (\delta s_{\rm min})\big)$ & $\til{\cO}\big(d_\sA^{14} d_\sB^{11}/(\delta^2 s_{\rm min}^3)\big)$~\cite{metger2023stateqipspace} \\ \hline
    Purified sample & $\til{\cO}\big(1/(\delta s_{\rm min}^2)\big)$ & $\til{\cO}\big(d_\sA d_\sB / (\delta^2 s_{\rm min}^2)\big)$ & N/A \\ \hline
    Mixed sample & $\til{\cO}\Big(\f{d_\sA}{\delta^2 s_{\rm min}^3} \big(\f{1}{\rho_{\rm min}^2} + \f{1}{\sigma_{\rm min}^2}\big)\Big)$ & $\til{\cO}\Big(\f{d_\sA}{\delta^4 s_{\rm min}^4} (r_\rho + r_\sigma)\Big)$ & N/A \\ 
\end{tabular}
\label{tab:compare Uhlmann and tomography}
\end{table*}
\renewcommand{\arraystretch}{1.0}









%=====================================================================================================================================================================================================




\subsection{Comparison with approaches based on state measurements}
\label{sec:comparison with state tom}


The most naive approach to implementing the Uhlmann transformation would be to use \emph{quantum state tomography}~\cite{ODonnell2016EffTomo, Haah2017samploptTomo, Guta2020FastTomo, Apeldoorn2023TomoPurifQuery, Chen2023WhenAdapTomo, hu2024sampleoptimalmemoryefficient}.
Using quantum state tomography, information about the states is obtained as classical data.
Then, the Uhlmann unitary in Eq.~\eqref{inteq:57} is directly computed.
By a brute-force decomposition of the computed unitary into one- and two-qubit gates~\cite{nielsen2010quantum}, a quantum circuit for the Uhlmann transformation can be obtained.
However, this approach clearly requires exponential query/sample complexity, as well as exponential classical time and space complexities in the number of bits, since it needs to handle matrices of exponential size in classical computation.
In addition, the number of gates arising from the decomposition is, in general, exponentially large.


Recently, an explicit quantum algorithm implementing the Uhlmann transformation was proposed in Ref.~\cite{metger2023stateqipspace}. This algorithm resolves the issue of exponential space use by sequentially encoding the measurement outcomes into small unitaries.
This approach enables the implementation of the Uhlmann transformation in polynomial space, while exponential time in the number of bits in classical computation and exponential query/sample complexity, $\poly(d_\sA d_\sB, 1/\delta)$, cannot be avoided, as it still relies on an exponential number of measurements.



Our algorithms do not require such a measurement on exponentially many copies of the states, and thus offer advantages in the query/sample complexity. 
Moreover, our algorithms only use a polynomial number of qubits at any one time.
Since our algorithms do not use classical data obtained by state measurements, they also avoid the exponential time and space complexities in classical computation that the naive tomography-based approach requires.



To compare our algorithms with the above approaches,
we summarize the query/sample complexity of these approaches in \Cref{tab:compare Uhlmann and tomography}.
For simplicity, we include only the results that have better scaling with respect to $\delta$ in this table.






In the purified query access model, our Uhlmann transformation algorithm clearly has an advantage over the other approaches.
We achieve an exponential speedup: our algorithm removes a factor of $d_\sA d_\sB$ in the query complexity and makes the dependence on $\delta$ logarithmic, demonstrating a significant improvement.
In the purified sample access model, our algorithm also remarkably outperforms the other approaches.

In the mixed sample access model, while our algorithm has advantages in terms of the dependence on $\delta$ and $s_{\rm min}$, the state tomography-based approach does not involve $\rho_{\rm min}$ and $\sigma_{\rm min}$.
Which algorithm achieves lower sample complexity depends on these values.
For instance, when these parameters satisfy $\delta^2 \leq \f{r_\rho + r_\sigma}{(1/\rho_{\rm min}^2 + 1/\sigma_{\rm min}^2)s_{\rm min}}$, our algorithm has an advantage.
In particular, when $r_\rho \approx 1/\rho_{\rm min}^2$ and $r_\sigma \approx 1/\sigma_{\rm min}^2$, for any $\delta$ such that $\delta \leq 1/\sqrt{s_{\rm min}}$, ours achieves a lower sample complexity than the naive state tomography-based approach.
We also stress that, even when our algorithms come with larger sample complexity than the tomography-based approach, our algorithms still have a significant advantage in the time and space complexities required for classical computation.


In \Cref{sec:Comparing with a naive tomography-based}, we provide detailed derivations of the computational costs by the naive state tomography-based approach, employing the most efficient known algorithms for quantum state tomography: $\tilde{\mathcal{O}}(kd/\epsilon)$ queries~\cite{Apeldoorn2023TomoPurifQuery} or $\tilde{\mathcal{O}}(kd/\epsilon^2)$ samples~\cite{ODonnell2016EffTomo, Haah2017samploptTomo} for a rank-$k$ state on a $d$-dimensional system within trace distance error $\epsilon$.
Here, the time and space complexities in classical computation required by this approach are tolerated. 
In addition, in the sample access model, it is allowed to perform collective measurements on multiple copies, possibly even an exponential number of them.

In \Cref{sec:metger and yuen's algorithm}, we provide an overview and performance evaluation of the approach proposed in Ref.~\cite{metger2023stateqipspace}, adapted to our setting, specifically, the purified query access model.
It is not immediately clear whether the approach of Ref.~\cite{metger2023stateqipspace} can be applied to the sample access models.
Even if it is applicable (with slight modifications), the resulting sample complexity would necessarily exceed the query complexity, since the purified query access model is computationally more powerful than the sample access models.







\subsection{Lower bound on the query and sample complexities for the Uhlmann transformation}
\label{sec:overview of lower bound}

Our explicit algorithms provide upper bounds on the computational complexity for the Uhlman transformation.
Here, we ask the converse question---what is the limit that can never be surpassed by any potential algorithm? 
To this end, we provide a lower bound on the query and sample complexities for the Uhlmann transformation.

The Uhlmann transformation can be applied to the task of fidelity estimation. Through fidelity estimation, we can certify the degree of mixedness of a given quantum state, a task known as the mixedness testing~\cite{ODonnell2021SpectrumTesting, MdW16surveyquantproptest, Wright2016howtolearnQS}.
Based on a result from Ref.~\cite{Liu2025EstimationTrace} on the mixedness testing in the purified query access model, we can derive a lower bound on the query complexity of fidelity estimation, which directly implies a lower bound for the Uhlmann transformation.
Consequently, we obtain the following corollary, with further details provided in \Cref{sec:lower bound}.



\begin{restatable}[Lower bound on the query complexity for the Uhlamnn transformation]{corollary}{corloweruhlquery}
\label{cor:lower Uhlmann query}
In the purified query access model, for $\delta \in (0, 1/9)$, there exists a pair of rank-$k$ states $\rho^\sA$ and $\sigma^\sA$ such that any quantum query algorithm that realizes a quantum channel $\cT^{\sB}$ satisfying
\begin{align}
    \rF(\cT^\sB(\ketbra{\rho}{\rho}^{\sA\sB}), \ket{\sigma}^{\sA\sB}) \geq \rF(\rho^\sA, \sigma^\sA) - \delta,
\end{align}
requires a total of $\Omega\big(k^{1/3}\big)$ queries to $U_\rho^{\sA\sB}$, $U_\sigma^{\sA\sB}$, and their inverses.
\end{restatable}


Note that the lower bound on the query complexity in the purified query access model also provides a lower bound on the sample complexity in the purified and mixed sample access models. 
This is because the purified query access model has the most powerful computational power among these models.


We compare the lower bound in \Cref{cor:lower Uhlmann query} with the upper bound obtained in \Cref{infthm:uhlmenn purif query}.
Since $r_{\rm min} = k$ for two rank-$k$ states $\rho^\sA$ and $\sigma^\sA$, the lower bound for the Uhlmann transformation is given by $\Omega(r_{\rm min}^{1/3})$ queries.
Thus, comparing this lower bound with the upper bound $\cO(r_{\rm min})$ queries in \Cref{infthm:uhlmenn purif query}, where $\delta$ is assumed to be a constant, we find a cubic gap between the upper and lower bounds with respect to $r_{\rm min}$.
Note that $r \leq r_{\rm min}$.
Closing this gap is an interesting open problem for future research.


As a brief comment from a complexity-theoretic perspective, testing the closeness of two quantum states is generally known to be $\mathsf{QSZK}$-complete, given access to unitaries that prepare their purifications~\cite{watrous2002limitQSZK, watrous2009QSZKquantumattacks}.
This implies that if a quantum algorithm that computes fidelity with polynomial time complexity exists, then $\mathsf{BQP} = \mathsf{QSZK}$~\cite{Rethinasamy2023estimatingdist}.
Since the Uhlmann transformation applies to fidelity estimation, it follows that unless $\mathsf{BQP} = \mathsf{QSZK}$, there is no polynomial time quantum algorithm implementing the Uhlmann transformation in general, which is consistent with prior suggestions from the context of the information recovery problem~\cite{harlow2013CompvsFirewall, aaronson2016moneyblackhole, brakerski2023blackholeradiation}.





%=====================================================================================================================================================================================================




\subsection{Applications}
\label{sec:overview application}


To demonstrate the usefulness of the Uhlmann transformation algorithm, we present three concrete applications: estimation of the square root fidelity, the decoupling approach~\cite{hayden2008decoupling, dupuis2010decoupling, Szehr2013decouple2design, dupuis2014one, preskill2016quantumshannon}, and algorithmic implementation of the Petz recovery map~\cite{petz1986sufficient, petz1988sufficiency}.
Comprehensive discussions for each application are provided in Secs.~\ref{sec:fidelity estimation}, \ref{sec:decoupling and uhlmann}, and \ref{sec:sample petz}, respectively.






%====================================================================================================




\subsubsection{Square root fidelity estimation}
\label{sec:overview of fidelity est}


The Uhlmann transformation algorithm is applicable to estimating the square root fidelity between states $\rho^\sA$ and $\sigma^\sA$.
The key idea of our approach is to first implement the transformation, and then run an optimal algorithm that estimates the absolute value of the inner product between two pure states~\cite{Wang2024OptimalFidelityPure, wang2024sampleoptimal}.
Due to the Uhlmann's theorem, the output coincides with the square root fidelity $\sqrt{\rF}(\rho^\sA, \sigma^\sA)$.
This approach provides algorithms for estimating the fidelity based on the Uhlmann transformation, in the purified query, purified sample, and mixed sample access models.
%in the purified query model



In the sample access models, our discussion so far has focused only on algorithms that use quantum states sequentially, one by one, \emph{non-collectively}.
Nevertheless, it is also possible to consider algorithms that manipulate multiple quantum states \emph{collectively}.
Collective algorithms potentially lead to a notable reduction in the number of samples, while they require sufficient space to store all states simultaneously. In contrast, non-collective algorithms may require more samples, while they allow us to reuse the qubits that hold the sampled states, thereby reducing the space cost.


Our results on the computational cost of the square root fidelity estimation in all of these settings are provided in the following theorem. 
An in-depth discussion is provided in~\Cref{sec:fidelity estimation}





%=====================================================================================================================================================================================================





\begin{table*}[t]
\centering
\caption{A comparison of the query and sample complexities for square root fidelity estimation within error $\delta$ between our algorithms and the prior best algorithms. For notation, see \Cref{tab:technical notation}. The left side summarizes results focusing on the scaling in $\kappa_\rho$ and $\kappa_\sigma$, the reciprocals of the minimum non-zero eigenvalues of $\rho^\sA$ and $\sigma^\sA$, respectively, while the right side focuses on their ranks.
The mark $\ddag$ indicates the results obtained by using Eq.~\eqref{inteq:135} under the assumption that either $\supp[\rho^\sA] \subset \supp[\sigma^\sA]$ or $\supp[\sigma^\sA] \subset \supp[\rho^\sA]$, and the mark $\S$ indicates the results obtained from the mixed sample access model in Ref.~\cite{gilyen2022improvedfidelity} by using the fact that $r_{\rm min} \leq \kappa_{\rm min}$.
To our knowledge, no previous work has directly addressed fidelity estimation in the purified sample access model, and thus the mark $\P$ indicates the results from the mixed sample access model, relying on the fact that the purified sample access model is computationally stronger than the mixed sample access model. The mark $\|$ indicates that the algorithms leading to the results perform collective operations over multiple samples of the states.}
\renewcommand{\arraystretch}{2.0}
\begin{tabular}{>{\centering\arraybackslash}m{7em}|
                >{\centering\arraybackslash}m{9em}
                >{\centering\arraybackslash}m{10em}
                >{\centering\arraybackslash}m{5em}
                >{\centering\arraybackslash}m{6em}}
                
& \parbox[c][2em][c]{19em}{\centering In terms of $\kappa$-scaling} 
& & \parbox[c][2em][c]{14em}{\hspace{-4mm}\centering In terms of $r$-scaling} & \\
\hline 
& This work &\hspace{-1mm} Previous work &\hspace{2mm}This work& Previous work \\
\hline
    \parbox[c][4em][c]{7em}{\centering Purified query}  
    & \parbox[c][4em][c]{9em}{\centering $^\ddag\til{\cO}\big(\f{\sqrt{\kappa_\rho\kappa_\sigma}}{\delta}\big)$} 
    & \hspace{-1mm}\parbox[c][4em][c]{9em}{\centering $^\ddag\til{\cO}\big(\f{\kappa_{\rm min}^2\kappa_\rho\kappa_\sigma}{\delta}\big)$~\cite{Liu2024geometricmean}}
    & \parbox[c][4em][c]{6em}{\centering $\til{\cO}\big(\f{r_{\rm min}}{\delta^2}\big)$}
    & \parbox[c][4em][c]{6em}{\centering $\til{\cO}\big(\f{r_{\rm min}^{2.5}}{\delta^5}\big)$~\cite{gilyen2022improvedfidelity}}
    \\ \hline
    \parbox[c][5em][c]{7em}{\centering Purified sample}
    & \parbox[c][5em][c]{9em}{\centering \makecell{$^{\|, \ddag}\til{\cO}\big(\f{\kappa_\rho\kappa_\sigma}{\delta^2}\big)$  \\[0.9ex]$^\ddag\til{\cO}\big(\f{\kappa_\rho\kappa_\sigma}{\delta^3}\big)$}} 
    & \parbox[c][5em][c]{9em}{\centering \makecell[c]{$^{\P, \S}\til{\cO}\big(\f{\kappa_{\rm min}^{5.5}}{\delta^{12}}\big)$~\cite{gilyen2022improvedfidelity} \\
    [0.9ex]\hspace{-3mm}$^{\P, \ddag}\til{\cO}\big(\f{\kappa_{\rm min}^5\kappa_\rho^2\kappa_\sigma^2}{\delta^3}\big)$~\cite{Liu2024geometricmean}}}
    & \parbox[c][5em][c]{6em}{\centering \makecell{$^\|\til{\cO}\big(\f{r_{\rm min}^2}{\delta^4}\big)$ \\
    [0.9ex]$\til{\cO}\big(\f{r_{\rm min}^2}{\delta^5}\big)$}}
    & \hspace{-3mm}\parbox[c][5em][c]{6em}{\centering $^\P\til{\cO}\big(\f{r_{\rm min}^{5.5}}{\delta^{12}}\big)$~\cite{gilyen2022improvedfidelity}}
    \\ \hline
    \parbox[c][5em][c]{7em}{\centering Mixed sample} 
    & \parbox[c][5em][c]{9em}{\centering \makecell[c]{$^{\|, \ddag}\til{\cO}\big(\f{\kappa_\rho\kappa_\sigma}{\delta^2}\big)$ \\
    [0.9ex]$^\ddag\til{\cO}\big(\f{d_\sA (\kappa_\rho\kappa_\sigma)^{3/2}(\kappa_\rho^2 + \kappa_\sigma^2)}{\delta^4}\big)$}}
    & \parbox[c][5em][c]{9em}{\centering \makecell[c]{$^\S\til{\cO}\big(\f{\kappa_{\rm min}^{5.5}}{\delta^{12}}\big)$~\cite{gilyen2022improvedfidelity} \\
    [0.9ex]\hspace{-1mm}$^\ddag\til{\cO}\big(\f{\kappa_{\rm min}^5\kappa_\rho^2\kappa_\sigma^2}{\delta^3}\big)$~\cite{Liu2024geometricmean}}}
    & \parbox[c][5em][c]{6em}{\centering \makecell[c]{$^\|\til{\cO}\big(\f{r_{\rm min}^2}{\delta^4}\big)$ \\
    [0.9ex]\hspace{1mm}$\til{\cO}\big(\f{d_\sA^3 r_{\rm min}^7}{\delta^{15}}\big)$}}
    & \parbox[c][5em][c]{6em}{\centering $\til{\cO}\big(\f{r_{\rm min}^{5.5}}{\delta^{12}}\big)$~\cite{gilyen2022improvedfidelity}}
    \\ 
\end{tabular}
\label{tab:fidelity}
\end{table*}
\renewcommand{\arraystretch}{1.0}









%=====================================================================================================================================================================================================





\begin{theorem}[Square root fidelity estimation using the Uhlmann transformation algorithms]
\label{thm:informal fidelity result}
In each of the three query/sample access models, for any $\delta \in (0, 1)$, there exists a quantum query/sample algorithm that outputs $\sqrt{\tilde{\rF}}$ satisfying $\big|\sqrt{\tilde{\rF}} - \sqrt{\rF}(\rho^\sA, \sigma^\sA)\big| \leq \delta$ with probability at least $2/3$. 
The computational costs of these algorithms in each model are as follows.
\begin{itemize}
\item The purified query access model: $q = \til{\cO}\big(\f{1}{\delta}{\rm min}\big\{\f{1}{s_{\rm min}}, \f{r}{\delta}\big\}\big)$ queries to $U_\rho^{\sA\sB}$, $U_\sigma^{\sA\sB}$, and their inverses. The quantum circuit consists of $\til{\cO}(q)$ one- and two-qubit gates, and $\cO(\log{(d_\sA d_\sB)} + \log{(1/\delta}))$ qubits suffice at any one time.
\item The purified sample access model: 
\begin{itemize}
    \item With collective operations: $n = \til{\cO}\big(\frac{1}{\delta^2} \min\big\{\f{1}{s_{\rm min}^2}, \f{r^2}{\delta^2}\big\}\big)$ samples of $\ket{\rho}^{\sA\sB}$ and $\ket{\sigma}^{\sA\sB}$. The quantum circuit consists of $\til{\cO}(n^4)$ one- and two-qubit gates, and $\til{\cO}(n^2)$ qubits suffice at any one time.
    \item With non-collective operations: $n = \til{\cO}\big(\frac{1}{\delta^3} \min\big\{\f{1}{s_{\rm min}^2}, \f{r^2}{\delta^2}\big\}\big)$ samples of $\ket{\rho}^{\sA\sB}$ and $\ket{\sigma}^{\sA\sB}$. The quantum circuit consists of $\til{\cO}(n)$ one- and two-qubit gates, and $\cO(\log{(d_\sA d_\sB)} + \log{(1/\delta)})$ qubits suffice at any one time.
\end{itemize}
\item The mixed sample access model: 
\begin{itemize}
    \item With collective operations: $n = \til{\cO}\big(\frac{1}{\delta^2} \min\big\{\f{1}{s_{\rm min}^2}, \f{r^2}{\delta^2}\big\}\big)$ samples of $\rho^\sA$ and $\sigma^\sA$. 
    The quantum circuit consists of $\til{\cO}(n^4)$ one- and two-qubit gates, and $\til{\cO}(n^2)$ qubits suffice at any one time.
    \item With non-collective operations: $n = \til{\cO}\Big(\f{d_\sA}{\delta^4}\min\Big\{\f{\kappa_\rho^2 + \kappa_\sigma^2}{s_{\rm min}^3}, \f{d_\sA^2 r^7}{\delta^{11}}\Big\}\Big)$ samples of $\rho^\sA$ and $\sigma^\sA$. The quantum circuit consists of $\til{\cO}(n)$ one- and two-qubit gates, and $\cO(\log{d_\sA} + \log{(1/\delta)})$ qubits suffice at any one time.
\end{itemize}
\end{itemize}

\end{theorem}



We summarize our results on the query and sample complexities in \Cref{tab:fidelity}, along with the most efficient results in prior works to our knowledge.
In comparison, we should note that $r \leq r_{\rm min} \leq \kappa_{\rm min}$. Moreover, as derived in \Cref{sec:appendix product of singvalue}, if either $\supp[\rho^\sA] \subset \supp[\sigma^\sA]$ or $\supp[\sigma^\sA] \subset \supp[\rho^\sA]$, we have
\begin{align}
\label{inteq:135}
    1/s_{\rm min} 
    &\leq 1/\sqrt{\rho_{\rm min}\sigma_{\rm min}} = \sqrt{\kappa_\rho\kappa_\sigma}, 
\end{align}
where $\supp[A]$ is the support of a matrix $A$, i.e., the orthogonal complement of the kernel of $A$.
By applying Eq.~\eqref{inteq:135} to the computational costs in \Cref{thm:informal fidelity result} for clarity, the results in \Cref{tab:fidelity} (marked by $\ddag$) are obtained.


In the purified query access model, our result achieves a significant improvement over the previous best-known results in Refs.~\cite{gilyen2022improvedfidelity, Liu2024geometricmean}.
The result in Ref.~\cite{Liu2024geometricmean} can be derived under the assumption that $\supp[\rho^\sA] \subset \supp[\sigma^\sA]$ or $\supp[\sigma^\sA] \subset \supp[\rho^\sA]$.
For simple comparison, we suppose $\kappa$ such that $\kappa_\rho, \kappa_\sigma \leq \kappa$.
While the result in Ref.~\cite{Liu2024geometricmean} scales as $\til{\cO}(\kappa^4/\delta)$, our result scales linearly as $\til{\cO}(\kappa/\delta)$, which demonstrates a quartic improvement.
Furthermore, the result in Ref.~\cite{gilyen2022improvedfidelity} scales as $\til{\cO}\big(r_{\rm min}^{2.5}/\delta^5\big)$, whereas ours scales as $\til{\cO}(r_{\rm min}/\delta^2)$, highlighting two improvements: linear scaling in $r_{\rm min}$ and a substantial reduction in the dependence on $\delta$.
We also demonstrate significant improvements in the purified sample access model, although this may be due to the model not having been thoroughly explored.




In the mixed sample access model, our results $\til{\cO}(\kappa^2/\delta^2)$ and $\til{\cO}(r_{\rm min}^2/\delta^4)$ with collective operations (marked by $\|$ in \Cref{tab:fidelity}) show notable improvement on the sample complexity over the previous best-known results $\til{\cO}(\kappa^9/\delta^3)$~\cite{Liu2024geometricmean} and $\til{\cO}(r_{\rm min}^{5.5}/\delta^{12})$~\cite{gilyen2022improvedfidelity}.
We achieve this by employing a recently developed technique known as the \emph{random purification}~\cite{tang2025conjugatequerieshelp}.
The random purification enables us to achieve the task in the mixed sample access model with roughly the square of the number of queries in the purified query access model, while it requires collective operations. 
By applying the random purification to the fidelity estimation, we highlight its remarkable power.
In \cite{tang2025conjugatequerieshelp}, the technique of random purification was used to argue the limitation of the power of access to state purification unitaries, which merely shows a quadratic improvement over access to sample copies.
On the other hand, our result shows that one can rather use this technique positively for designing an efficient sample-based protocol.


It is also noteworthy that, focusing on the $\delta$-scaling, we find that our results, $\til{\cO}(\kappa/\delta)$ queries and $\til{\cO}(\kappa^2/\delta^2)$ samples in the purified query and mixed sample access models, are nearly optimal. This is because the lower bounds on the query and sample complexities for square root fidelity estimation are given by $\Omega(1/\delta)$ and $\Omega(1/\delta^2)$~\cite{Liu2024geometricmean}, respectively.
The optimalities with respect to the $\kappa$-scaling are still open.




We have not yet found an algorithm in the mixed sample access model that uses only sequential operations and shows an advantage in sample complexity. Although the Uhlmann transformation inherently requires purified states, no algorithm is currently known that realizes purification solely through sequential operations, with a number of samples that scales with the rank rather than the dimension.
If such a purification algorithm were developed, the Uhlmann transformation would become more tractable. This could improve the sample complexity for fidelity estimation in the mixed sample access model, even without collective operations.






We comment on the number of qubits in the cases with and without collective operations.
When we do not use collective operations, $\cO\big(\log{(d_\sA d_\sB)} + \log{(1/\delta)}\big)$ qubits suffice at any time in the algorithms. The term $\cO(\log(d_\sA d_\sB))$ comes from our Uhlmann transformation algorithms, while $\cO(\log{(1/\delta)})$ is from the algorithm in Ref.~\cite{Wang2024OptimalFidelityPure, wang2024sampleoptimal}, which is based on quantum phase estimation and is employed in our fidelity estimation algorithms.
On the other hand, when we employ collective operations, $\til{\cO}\big(n^2)$ qubits are used simultaneously, where $n$ is the number of sampled states. 
This is due to the Schur transform applied in the random purification~\cite{tang2025conjugatequerieshelp} and follows from the results in Ref.~\cite{burchardt2025highdimensionalschur}.
As Theorem~\ref{thm:informal fidelity result} shows, collective operations can substantially reduce the sample cost $n$; however, the number of qubits required at once increases with $n$. 
Hence, a non-collective approach may be more suitable when the number of simultaneously available qubits is limited.













\subsubsection{Information-theoretic tasks with the decoupling approach}
\label{sec:overview of decoup}


To accomplish the theoretical limits of various information-theoretic tasks, the decoupling approach~\cite{hayden2008decoupling, dupuis2010decoupling, Szehr2013decouple2design, dupuis2014one, preskill2016quantumshannon}, which relies on the existence of the Uhlmann transformation, is widely employed (see e.g., Refs.~\cite{hayden2008decoupling, Abeyesinghe2009Motherfamilytree, Datta2011apexfamilytree}).
Although the existence of the Uhlmann transformation is guaranteed, its implementation methods have long been of interest in the context of the decoupling approach.

Within the decoupling approach, we apply the proposed Uhlmann transformation algorithms to information-theoretic tasks and evaluate the computational costs. 
We particularly focus on two tasks: \emph{entanglement transmission}~\cite{Schumacher1996sendingentanglement, schumacher2001approximateerrorcorrection, hayden2007black, datta2013oneshotentassistQCcom, beigi2016decoding, khatri2024principlemodern} and \emph{quantum state merging}~\cite{Horodecki2005Partialquantinfo, berta2008singleshot, dupuis2014one, yamasaki2019statemergesmalldimension}.
Comprehensive discussions, including the settings, are provided in \Cref{sec:decoupling and uhlmann}.
The main results can be found in \Cref{thm:Uhlmann decoder ent trans} for entanglement transmission and in \Cref{thm:Uhlmann state merging} for quantum state merging.
Our results bring us a step closer to achieving the theoretical limits of these tasks, providing explicit methods toward their realization.





\subsubsection{Algorithmic implementation of the Petz recovery map}
\label{sec:overview of petz}


The Petz recovery map~\cite{petz1986sufficient, petz1988sufficiency} is one of the few powerful tools for reversing the effect of a noisy quantum channel, achieving nearly optimal performance and having an explicit form.
Quantum algorithms for the Petz recovery map are studied, for example, in Ref.~\cite{gilyen2022petzmap, biswas2023noiseadapted}, and have demonstrated significant progress toward its implementation.
These algorithms, however, require access to a Stinespring unitary of the noisy quantum channel, which can be a practically challenging requirement.
We introduce a method for directly implementing the Petz recovery map from a noisy quantum channel without relying on its Stinespring unitary.

The key idea is to approximate the Stinespring unitary using the Uhlmann transformation algorithm for the mixed sample access model. 
In fact, the Uhlmann unitary $U^\sS$ approximately mapping $\ket{\Phi}^{\sR\sA}\ket{0}^{\sF}$ to $\ket{\tau_{\rm c}}^{\sR\sB\hat{\sR}\hat{\sB}}$ coincides with one of the Stinespring unitaries of the noisy quantum channel $\cF^{\sA\rarr\sB}$, where $\sS=\sA\sF = \sB\hat{\sR}\hat{\sB}$.
The state $\ket{\Phi}^{\sR\sA}$ is a maximally entangled state and the state $\ket{\tau_{\rm c}}^{\sR\sB\hat{\sR}\hat{\sB}}$ is a canonical purification of the \emph{Choi--Jamio\l okwki state} of $\cF^{\sA\rarr\sB}$. 
We evaluate the number of uses of the quantum channel $\cF^{\sA\rarr\sB}$ for implementing the Petz recovery map.
A detailed discussion is in \Cref{sec:sample petz}.






%======================================================================================================================================================


